\subsection{Três liquidações para a mesma ordem}
\label{subsec:garantia-onze}

Imaginemos que, desta vez, a compra de Alice a Bob passa por um depósito
em garantia. O preço continua a ser 7. Já foram reconhecidas na carteira
comum 11 unidades, e uma unidade cobrirá a taxa da transacção que liquidar
a ordem. As 11 são o valor que entrou nessa carteira. A transacção que a
financiou teve as suas próprias entradas, eventuais saídas de troco e taxa.
Começamos a nossa conta depois desse financiamento.

Alice é a compradora \(B\), Bob é o vendedor \(V\) e o moderador é
\(M\). A ordem estabelece as condições de entrega, autoriza o reembolso
integral em caso de cancelamento e fixa em uma unidade o máximo a pagar
ao moderador se for precisa uma decisão. Dos 11 depositados estarão
disponíveis para os beneficiários \(11-1=10\). O que se discute entre
as partes é a distribuição desses dez.

\begin{center}
\begin{tabular}{l|rrr|r}
resultado & Alice & Bob & moderador & taxa da transacção\\
\hline
entrega aceite & 3 & 7 & 0 & 1\\
compra cancelada & 10 & 0 & 0 & 1\\
entrega parcial, com decisão & 4 & 5 & 1 & 1
\end{tabular}
\end{center}

Na primeira linha Alice recebeu o bem, confirma o preço e recupera o que
sobra, \(11-7-1=3\). A assinatura pode ser concluída por Alice e Bob.
Na segunda linha a compra foi cancelada por acordo; o comprador recebe
\(11-1=10\). Os mesmos dois signatários autorizam uma distribuição
inteiramente diferente. Na terceira linha o moderador decidiu que a
parte entregue vale 5. A ordem prevê para o seu serviço uma remuneração
de 1. Restam ao comprador \(11-5-1-1=4\). Uma das unidades paga o
serviço comercial da moderação e tem um endereço de recepção. A outra
paga o espaço e processamento da transacção e entra no campo da taxa.
Esta distinção fica especialmente clara quando os dois pagamentos têm
o mesmo valor.

As três linhas descrevem direitos económicos. Uma componente nula
significa que a ordem nada atribui àquela parte. O formato concreto
da transacção pode repartir a soma de um beneficiário por mais de uma
saída. O que o signatário confere é o total dirigido a cada destinatário
autorizado, bem como o troco e a taxa que a proposta declara. Para
executar o reembolso, por exemplo, o vector exige que as saídas
destinadas ao comprador totalizem dez. Uma divisão em mais de uma
saída conserva o mesmo resultado económico, desde que a transacção
satisfaça as regras aplicáveis.

\paragraph{O que chega ao nó que verifica a liquidação.}
Tomemos a primeira linha e suponhamos que uma única saída antiga
financiou o depósito. Se o pseudo-compromisso de entrada tem máscara
\(x\), o termo da entrada é \(xG+11H\). Sejam \(a\) e \(b\) as
máscaras das saídas de Alice e Bob, com \(a+b=x\) no corpo de
escalares. O balanço público expande-se como
\[
 (aG+3H)+(bG+7H)+1H=(a+b)G+(3+7+1)H=xG+11H.
\]
Se a ordem for liquidada segundo a terceira linha, haverá distribuições
de máscaras \(a',b',m'\) com \(a'+b'+m'=x\), tomando para a comparação
o mesmo valor de máscara do lado da entrada. A soma do lado direito da
liquidação passa a ser
\[
 (a'G+4H)+(b'G+5H)+(m'G+1H)+1H=xG+11H.
\]
A escolha de máscaras pode ser renovada ao construir cada proposta.
O exemplo fixa a sua soma para se ver à primeira leitura o que as duas
contas têm em comum. A igualdade de pontos atesta o equilíbrio em ambas,
e as provas de intervalo e de gasto completam as verificações ordinárias.
A mensagem comercial liga a distribuição correcta ao estado da entrega.
É por isso que se verifica a proposta antes de fornecer a contribuição
de assinatura.

Se, por fim, a taxa de liquidação subir de 1 para 2 e o preço se
mantiver, o primeiro resultado passará a \((2,7,0)\). A soma é agora 9,
e \(9+2=11\). O resultado antigo \((3,7,0)\), combinado com taxa 2,
exigiria 12 de uma entrada de 11. Uma simples mudança na data de envio
pode, portanto, obrigar os participantes a concordar numa nova proposta
e em quem suporta a diferença. Nem a decisão comercial nem a imagem
de chave da saída financiada mudam por causa dessa nova conta.
